\newpage
\section{Spécification du projet}
\label{sec:cdc}

\paragraph{} Nous avons présenté l'objectif du projet dans la section \ref{sec:intro}. 
Dans cette section, nous présentons la spécification complète de notre logiciel, correspondant 
au cahier des charges établi en début de projet. Nous y détaillons les notions fondamentales, 
les contraintes techniques, la modélisation de l'environnement, ainsi que l'ensemble des 
fonctionnalités attendues du système domotique intelligent.

% -----------------------------------------------------------------------
\subsection{Notions de base et contraintes du projet}
\label{sec:spec1}

\subsubsection{Fonctionnement général du logiciel}

Le logiciel simule une maison intelligente en vue verticale 2D, composée de plusieurs pièces 
(salon, chambre, cuisine, salle de bain) reliées par un couloir central. Un agent autonome, 
appelé \textit{Maître de maison}, évolue dans cet environnement selon une routine journalière 
prédéfinie. Le système domotique analyse en continu l'état du maître — ses besoins 
physiologiques et son humeur — afin de déclencher automatiquement des actions adaptées sur 
les équipements connectés de l'habitat.

\paragraph{} La simulation se déroule en temps accéléré et configurable : une journée complète 
de 24 heures simulées correspond à 24 minutes réelles, soit un rapport de 1 minute réelle pour 
1 heure simulée. Ce mode de fonctionnement permet d'observer l'ensemble du comportement du 
système sur une journée complète dans un laps de temps réduit, tout en conservant la cohérence 
temporelle des événements simulés.

\paragraph{} L'architecture logicielle repose sur une séparation claire des responsabilités : 
un moteur de simulation gère la logique métier (déplacements, besoins physiologiques, 
déclenchement des actions), une couche graphique assure l'affichage de l'environnement et 
de l'état du maître, et un module de configuration centralise les paramètres globaux de la 
simulation. Cette organisation en couches facilite la maintenabilité du code et l'évolution 
future du projet.

\subsubsection{Notions et terminologies de base}

\paragraph{Domotique} La domotique rassemble l'ensemble des techniques permettant de 
contrôler, programmer et automatiser une habitation. Elle regroupe les domaines de 
l'électronique, de l'informatique, de la télécommunication et des automatismes. Le terme 
est issu du latin \textit{domus} (domicile) et du suffixe \textit{-tique} (technique). Dans 
le cadre de ce projet, la domotique est simulée logiciellement : les équipements physiques 
sont remplacés par des objets Java dont l'état est géré et affiché en temps réel.

\paragraph{Maître de maison} Il s'agit de l'utilisateur principal du système domotique, 
pour lequel le système s'adapte et se personnalise. Il possède un profil complet incluant 
son état d'humeur ainsi que ses besoins physiologiques : faim, soif, fatigue et hygiène. 
Chacun de ces paramètres est modélisé par une barre de jauge évoluant de 0 à 100\%, dont 
la valeur influence directement les décisions du système. Un seul maître de maison est 
géré dans cette simulation.

\paragraph{État d'humeur} L'humeur représente l'état émotionnel du maître à un instant 
donné. Elle est définie aléatoirement au début de la simulation, c'est-à-dire au réveil 
du maître, et évolue dynamiquement en fonction des besoins physiologiques, de l'état de 
santé et des imprévus survenus au cours de la journée. Un état d'humeur négatif est 
engendré par une dégradation des paramètres physiologiques, tandis qu'une stabilité de 
ceux-ci favorise une humeur positive. Les quatre états d'humeur possibles sont présentés 
dans le tableau \ref{tab:humeurs}.

\begin{table}[h!]
\centering
\begin{tabular}{|p{3cm}|p{9cm}|}
\hline
\textbf{État d'humeur} & \textbf{Description et influence sur le système} \\
\hline
Joyeux & État positif. Le système favorise les activités conviviales et les repas 
colorés. L'ambiance lumineuse est vive et chaleureuse. \\
\hline
Triste & État négatif. Le système propose des activités réconfortantes, 
des plats doux et une ambiance apaisante. \\
\hline
En colère & État négatif. Le système limite les stimulations extérieures, 
propose des repas consistants et une musique calme. \\
\hline
Neutre & État stable. Le système propose de la nouveauté et de la variété 
pour stimuler la curiosité du maître. \\
\hline
\end{tabular}
\caption{États d'humeur du maître de maison et influence sur le système}
\label{tab:humeurs}
\end{table}

\paragraph{Besoins physiologiques} Les besoins du maître de maison sont modélisés sous 
forme de barres de jauge évoluant continuellement au cours de la simulation. Chacune 
évolue selon un taux horaire défini, comme présenté dans le tableau \ref{tab:besoins}.

\begin{table}[h!]
\centering
\begin{tabular}{|p{3cm}|p{4cm}|p{5cm}|}
\hline
\textbf{Besoin} & \textbf{Taux de diminution} & \textbf{Condition de remontée} \\
\hline
Faim & $-7\%$ par heure simulée & Après un repas \\
\hline
Soif & $-15\%$ par heure simulée & Après hydratation \\
\hline
Énergie & $-5$ à $-20\%$ par heure selon l'activité & Durant le sommeil \\
\hline
Hygiène & $-3\%$ par heure simulée & Après une douche ou un bain \\
\hline
\end{tabular}
\caption{Évolution des besoins physiologiques du maître de maison}
\label{tab:besoins}
\end{table}

\paragraph{Objets connectés} Plusieurs équipements domotiques sont simulés dans la maison. 
Chaque objet possède un état actif ou inactif géré automatiquement par le système :
\begin{itemize}
    \item \textbf{Éclairage} : les lumières de chaque pièce peuvent être allumées ou 
    éteintes automatiquement. Leur intensité (forte ou faible) ainsi que leur couleur 
    (LED colorées) sont également gérées en fonction de l'humeur et de l'heure.
    \item \textbf{Chauffage} : la maison régule automatiquement la température intérieure 
    afin de maintenir un confort ambiant, ou de l'adapter à l'humeur du maître.
    \item \textbf{Volets} : leur ouverture et fermeture sont automatisées selon la météo 
    ou l'horaire de la journée.
    \item \textbf{Enceinte connectée} : l'ambiance sonore est gérée selon plusieurs 
    modes : calme, dynamique ou relaxant, selon le contexte émotionnel du maître.
\end{itemize}

\paragraph{Routine journalière} Le maître de maison suit une routine prédéfinie qui varie 
selon le type de journée. Les tableaux \ref{tab:semaine} et \ref{tab:weekend} présentent 
respectivement le déroulement d'une journée de semaine et d'une journée de weekend.

\begin{table}[h!]
\centering
\begin{tabular}{|p{2.5cm}|p{9.5cm}|}
\hline
\textbf{Horaire} & \textbf{Activité} \\
\hline
07:00 & Réveil — ouverture automatique des fenêtres. Barre de fatigue au maximum, 
barre d'hygiène basse. Le système propose une douche. \\
\hline
07:05 & Douche (20 min) — barre d'hygiène au maximum. Barre de faim basse. 
Le système propose un petit-déjeuner. \\
\hline
07:25 & Petit-déjeuner — barre de faim au maximum. \\
\hline
07:40 & Départ au travail — la porte de sortie s'ouvre. \\
\hline
18:40 & Retour à domicile — humeur basse. Le système propose une douche, 
un repas, une boisson et une activité relaxante. \\
\hline
20:00 & Dîner. \\
\hline
22:00 & Coucher — fin de journée. \\
\hline
\end{tabular}
\caption{Routine journalière en semaine}
\label{tab:semaine}
\end{table}

\begin{table}[h!]
\centering
\begin{tabular}{|p{2.5cm}|p{9.5cm}|}
\hline
\textbf{Horaire} & \textbf{Activité} \\
\hline
10:00 & Réveil tardif — proposition de douche. \\
\hline
10:30 & Petit-déjeuner. \\
\hline
11:00 & Travail à domicile dans le bureau. \\
\hline
14:00 & Déjeuner. \\
\hline
14:30--15:30 & Sieste — légère hausse de la barre d'énergie. \\
\hline
16:00--18:00 & Sortie extérieure — diminution des barres de fatigue, d'hygiène 
et d'humeur. \\
\hline
18:00 & Retour — activité relaxante proposée par le système. \\
\hline
21:00 & Dîner puis coucher à 22:00. \\
\hline
\end{tabular}
\caption{Routine journalière le weekend}
\label{tab:weekend}
\end{table}

\paragraph{Imprévus} Des événements non planifiés peuvent survenir au cours de la journée 
et perturber la routine normale du maître. Le système doit détecter ces situations et 
réagir de manière adaptée. Les imprévus pris en compte dans la simulation sont les suivants :
\begin{itemize}
    \item \textbf{Réveil en retard} : entraîne le décalage de l'ensemble des activités 
    matinales et une dégradation plus rapide des besoins physiologiques.
    \item \textbf{Message téléphonique urgent} : sollicitation professionnelle imprévue 
    en dehors des horaires habituels, nécessitant une réorganisation des priorités.
    \item \textbf{Retour tardif du travail} : décalage du dîner, annulation ou report 
    des activités du soir.
\end{itemize}
Chaque imprévu influence négativement l'humeur du maître et modifie dynamiquement 
les priorités d'action du système.

\subsubsection{Contraintes et limitations connues}

\paragraph{}

\paragraph{} :


\paragraph{} Outils de développement utilisés :
\begin{enumerate}
    \item Langage de programmation : Java
    \item Environnement de développement : Eclipse
    \item Outil de rédaction du rapport : LATEX
\end{enumerate}

% -----------------------------------------------------------------------
\subsection{Fonctionnalités attendues du projet}
\label{sec:spec2}

\paragraph{} Cette sous-section présente l'ensemble des fonctionnalités que le système 
domotique doit implémenter. Celles-ci sont regroupées en deux catégories : les 
fonctionnalités générales de simulation, et les actions automatiques déclenchées en 
fonction de l'état du maître.

\paragraph{Fonctionnalités générales} Le système domotique doit assurer les 
fonctionnalités suivantes :
\begin{itemize}
    \item Modéliser graphiquement un appartement en 2D composé de plusieurs pièces 
    distinctes, d'un couloir central et d'une sortie.
    \item Simuler le déplacement autonome du maître de maison dans l'appartement, 
    en respectant les contraintes physiques telles que les murs, les portes et les 
    bordures de la carte.
    \item Gérer en continu l'évolution des besoins physiologiques du maître selon 
    les taux horaires définis dans le tableau \ref{tab:besoins}.
    \item Calculer dynamiquement l'état d'humeur du maître à partir de ses paramètres 
    physiologiques et des événements survenus.
    \item Gérer la routine journalière du maître selon le type de journée (semaine 
    ou weekend), en déclenchant les activités aux horaires prévus.
    \item Détecter et gérer les imprévus en notifiant le maître et en réorganisant 
    les priorités du système en conséquence.
    \item Permettre à l'utilisateur de mettre en pause, reprendre ou accélérer 
    la simulation via l'interface graphique.
    \item Afficher en temps réel l'état du maître (barres de besoins, humeur, 
    heure simulée) ainsi que l'état des équipements connectés.
\end{itemize}

\paragraph{Actions automatiques selon l'état du maître} Le système surveille en 
permanence l'état du maître et déclenche des actions spécifiques dès qu'un seuil 
prédéfini est atteint. L'ensemble de ces déclencheurs et de leurs effets est 
récapitulé dans le tableau \ref{tab:actions}.

\begin{table}[h!]
\centering
\begin{tabular}{|p{4cm}|p{4cm}|p{5cm}|}
\hline
\textbf{Condition déclenchante} & \textbf{Action du système} & \textbf{Effet} \\
\hline
Barre de faim $<$ 40\% aux heures de repas & Proposition d'un repas adapté 
à l'humeur & Remplissement de la barre de faim, légère hausse de l'humeur \\
\hline
Barre de soif $<$ 50\% & Proposition d'une boisson adaptée à l'humeur & 
Remplissement de la barre de soif \\
\hline
Barre d'énergie $<$ 20\% ou heure $>$ 22h30 & Activation du mode nuit 
(lumières tamisées, volets fermés, musique douce) & Remplissement de 
la barre de sommeil \\
\hline
Niveau d'hygiène $<$ 20\% & Proposition d'un bain immédiat & Hausse 
du niveau d'hygiène et effet positif sur l'humeur \\
\hline
Humeur négative (tristesse ou colère) & Ajustement de l'ambiance intérieure 
(éclairage, température, musique) & Retour progressif à un état 
émotionnel stable \\
\hline
Survenue d'un imprévu & Notification du maître et réorganisation 
des priorités & Limitation de l'impact sur le confort et le bien-être \\
\hline
\end{tabular}
\caption{Actions automatiques du système selon l'état du maître}
\label{tab:actions}
\end{table}

\paragraph{} Comme illustré dans le tableau \ref{tab:actions}, les actions du système 
sont hiérarchisées selon un ordre de priorité strict. Les besoins vitaux faim, soif 
et santé  sont toujours traités en premier, suivis du sommeil, de l'hygiène, du 
bien-être émotionnel, et enfin des activités de loisir. Cette organisation garantit 
une simulation cohérente, réaliste et centrée sur le bien-être du maître de maison.

\paragraph{Adaptation des repas selon l'humeur} Le système ne se contente pas de 
proposer un repas lorsque la barre de faim est basse : il adapte également le type 
de repas à l'état émotionnel du maître. Ainsi :
\begin{itemize}
    \item \textbf{Joyeux} : repas coloré et convivial (pizza, tacos, burgers).
    \item \textbf{Triste} : plat réconfortant à texture douce (purée, lasagnes, 
    hachis parmentier), avec éventuellement un dessert sucré si l'humeur est 
    particulièrement basse.
    \item \textbf{En colère} : plat consistant et épicé (kebab, steak, wrap).
    \item \textbf{Neutre} : cuisine du monde pour stimuler la curiosité 
    (mexicain, libanais, asiatique).
\end{itemize}

\paragraph{Adaptation des boissons selon l'humeur} De la même manière, le choix de 
la boisson proposée est conditionné par l'état émotionnel du maître :
\begin{itemize}
    \item \textbf{Humeur stimulante (joyeux, en colère)} : boisson rafraîchissante 
    et pétillante (soda, eau aromatisée, smoothie aux fruits).
    \item \textbf{Humeur basse (triste, neutre)} : boisson réconfortante et 
    énergisante (thé vert, thé noir, smoothie exotique).
\end{itemize}

\paragraph{} L'ensemble de ces fonctionnalités constitue le cœur du système domotique 
intelligent et adaptatif développé dans ce projet. Leur implémentation repose sur 
une architecture orientée objet rigoureuse, détaillée dans la section suivante.